\section{Further directions}\label{sec:discussion}
The four-branch argument converts terminal fourth moments into weighted
second moments. For five branch vertices, the proof allocates central arms to different
moment estimates and introduces additional boundary quantities, among
them $I_q$. Each skeleton has a separate assembly. An extension
to more branch vertices requires estimates that remain valid for the
composite messages produced by further mergers.

Known attachment constructions already produce families with unbounded
branch count~\cite[Theorem 1.3 and Proposition 5.2]{HMNTW}. Attaching two
copies with the same internal orientation and rooted message $m\ge0$,
where $\E m\le1$, through oppositely directed attachment edges contributes
$(Am)(Tm)=(\E m)^2-(Km)^2\le1$, the kernel form of those constructions. Asymmetric composite branches require further boundary estimates.

Two branches give explicit losses, $e^{-d/2000}$ for the two-outlet
path construction and $e^{-d/160000}$ for the nonsingleton central
allocations. The other constructions yield $Z_D\le1$.
